\documentclass[11pt]{amsart}

\usepackage[T1]{fontenc}
\usepackage{lmodern}
\usepackage{microtype}
\usepackage{textcomp}
\usepackage{upquote}   % so the backticks in the graph6 strings below are ASCII 0x60,
                       % and survive a copy-paste out of the PDF
\usepackage{mathtools,amssymb,amsthm}
\usepackage[hidelinks]{hyperref}

\hypersetup{
  pdftitle={Corner-free 4-cop-win graphs on 20 vertices, and a retraction question
            of Turcotte and Yvon}
}

\newtheorem{theorem}{Theorem}[section]
\newtheorem{lemma}[theorem]{Lemma}
\newtheorem{corollary}[theorem]{Corollary}
\theoremstyle{remark}
\newtheorem*{remark}{Remark}

\newcommand{\cn}{c}

\title[Corner-free 4-cop-win graphs on 20 vertices]
{Corner-free 4-cop-win graphs on 20 vertices,\\ and a retraction question of Turcotte and Yvon}

\date{}

\subjclass[2020]{05C57, 05C69, 05C75}
\keywords{cops and robbers, cop number, retract, corner, girth, counterexample}

\begin{document}

\begin{abstract}
Let $M_k$ be the least order of a connected graph with cop number exactly $k$. Turcotte and
Yvon, who proved $M_4=19$, ask in the discussion of their paper whether the $k$-cop-win graphs
on $M_k+1$ vertices always retract onto $k$-cop-win graphs on $M_k$ vertices, and single out
the case $k=4$. We answer that case negatively: two connected $4$-regular graphs of girth~$5$
on $20=M_4+1$ vertices have cop number exactly $4$ and have \emph{no corner}, hence no retract
of order $19$ whatsoever. They are exhibited below by their graph6 strings and their adjacency
lists, and the certificate is short enough to check by hand: girth at least $5$ with $\delta=4$
gives $\cn\ge4$ by a theorem of Aigner and Fromme, an explicit three-element set whose
complement induces a forest gives $\cn\le4$, and girth at least $5$ with $\delta\ge2$ forbids
corners. Only $k=4$ is settled here: $k=2$ is already false in the source, $k=3$ is true there
by its own classification, and $k\ge5$ is untouched because $M_5$ is unknown.
\end{abstract}

\maketitle

\section{The question}\label{sec:question}

We use the conventions of Turcotte and Yvon \cite{TY}. In the game of Cops and Robbers on a
finite connected simple graph $G$, the cops first choose their vertices, then the robber
chooses his; thereafter the cops move, then the robber, alternately, and on his turn a player
may stay put or move to an adjacent vertex. The cops win if a cop ever occupies the robber's
vertex, and $\cn(G)$ is the least number of cops with a winning strategy. Following
\cite{TY}, $G$ is \emph{$k$-cop-win} if $\cn(G)=k$ \emph{exactly}, and
\[
  M_k=\min\{|V(G)| : G \text{ connected and } \cn(G)=k\}.
\]
Thus $M_2=4$ (attained by $C_4$), $M_3=10$ (attained by the Petersen graph), and the main
theorem of \cite{TY} is $M_4=19$, attained by the Robertson graph \cite{Robertson}. A graph $H$
is a \emph{retract} of $G$ if $H$ is an \emph{induced} subgraph of $G$ and there is a map
$f\colon V(G)\to V(H)$ with $f|_{V(H)}=\mathrm{id}$ such that $xy\in E(G)$ implies
$f(x)f(y)\in E(H)$ or $f(x)=f(y)$.

In the discussion following their classification of the $3$-cop-win graphs on $11$ vertices,
the authors ask the following. We quote it in full, from the arXiv e-print
\texttt{2006.02998v2} of \cite{TY}.

\begin{quote}
``This is an interesting phenomenon: the unique $3$-cop-win graph on $10$ vertices is a
retract of all the $3$-cop-win graphs on $11$ vertices. This behaviour does not occur for the
$2$-cop-win graphs: the minimum $2$-cop-win graph is the $4$-cycle, on which the $5$-cycle
does not retract. Although we will not have any answer for this question in this article, it
would be interesting to know whether in general (even for $4$-cop-win graphs only), the
$k$-cop-win graphs on $M_k+1$ vertices can be retracted on $k$-cop-win graph(s) on $M_k$
vertices.''
\end{quote}

This is an unnumbered question in running prose, and not the numbered Conjecture of \cite{TY},
which is a different statement -- that no connected graph of order $19$ with maximum degree $5$
or $6$ has cop number $4$ -- and is not addressed here; in particular we do not claim, and
\cite{TY} does not prove, that the Robertson graph is the only $4$-cop-win graph on $19$
vertices. The authors state that $k=2$ already fails, their witness being $C_5$, which does not
retract onto $C_4$; and $k=3$ is true, since the classification just cited shows that all six
$3$-cop-win graphs on $11$ vertices retract onto the Petersen graph. We settle the case the
authors name, $k=4$, in the negative. Since $M_4=19$, the relevant order is $M_4+1=20$.

\section{The two graphs}\label{sec:graphs}

Let $W_1$ and $W_2$ be the graphs on the vertex set $\{0,1,\dots,19\}$ given by the following
graph6 strings, one $33$-byte line each, $33=1+\lceil\binom{20}{2}/6\rceil=1+32$:

\begin{verbatim}
W1 = S????B?K?W?WIg`cIP?j?QSA`GCoGBGG?
W2 = S??CA?_c?o@_Q_iOAcU@OK`@cO?MO?Bg?
\end{verbatim}

\noindent
Equivalently, by their adjacency lists (the label $v$ followed by the four neighbours of $v$):

{\footnotesize
\begin{verbatim}
W1                                          W2
 0:8,13,17,18   1:8,14,16,19                 0:6,9,13,15    1:7,12,15,17
 2:9,12,17,19   3:9,14,15,18                 2:8,13,16,17   3:9,10,14,16
 4:10,12,16,18  5:10,13,15,19                4:10,11,12,13  5:11,14,17,18
 6:11,12,13,14  7:11,15,16,17                6:0,12,16,18   7:1,13,18,19
 8:0,1,12,15    9:2,3,13,16                  8:2,14,15,19   9:0,3,17,19
10:4,5,14,17   11:6,7,18,19                 10:3,4,15,18   11:4,5,16,19
12:2,4,6,8     13:0,5,6,9                   12:1,4,6,14    13:0,2,4,7
14:1,3,6,10    15:3,5,7,8                   14:3,5,8,12    15:0,1,8,10
16:1,4,7,9     17:0,2,7,10                  16:2,3,6,11    17:1,2,5,9
18:0,3,4,11    19:1,2,5,11                  18:5,6,7,10    19:7,8,9,11
\end{verbatim}
}

Each is connected, $4$-regular on $20$ vertices with $40$ edges, and of girth $5$; the
$5$-cycles $0\,8\,15\,3\,18\,0$ in $W_1$ and $0\,6\,12\,1\,15\,0$ in $W_2$ show the girth is
not more than $5$, and the absence of triangles together with
$\max_{a\ne b}|N(a)\cap N(b)|=1$ shows it is not less.

Neither graph is new here: they are members of Meringer's census \cite{Meringer} of connected
$4$-regular graphs of girth at least $5$, being Graph~2 and Graph~1 of the file
\texttt{20\_4\_5.asc}, and their automorphism groups have orders $96$ and $20$ respectively,
which is how the identification is checked here.

\section{The certificate}\label{sec:cert}

Write $N[S]=S\cup\bigcup_{v\in S}N(v)$. A vertex $x$ is a \emph{corner} (or is
\emph{dominated}) if there is $u\ne x$ with $N(x)\subseteq N[u]$.

\begin{theorem}[Aigner and Fromme \cite{AF}]\label{thm:af}
If $G$ has girth at least $5$, then $\cn(G)\ge\delta(G)$.
\end{theorem}

\begin{lemma}\label{lem:forest}
Let $G$ be connected and suppose some $S\subseteq V(G)$ with $|S|=3$ has
$G[V(G)\setminus N[S]]$ a forest. Then $\cn(G)\le4$.
\end{lemma}

\begin{proof}
Place three cops on $S$ and let them never move. The robber cannot choose a vertex of $N[S]$,
since a cop would step onto him at once, and for the same reason he can never enter $N[S]$
later; so he is confined for the whole game to $R=V(G)\setminus N[S]$, and in fact to the
component $T$ of $G[R]$ containing his initial vertex. As $G[R]$ is a forest, $T$ is a tree.
A fourth cop walks to $T$ -- possible since $G$ is connected -- and then plays the one-cop
winning strategy on the tree $T$, in which the robber's legal moves are exactly the edges of
$T$. Hence four cops suffice.
\end{proof}

This argument is not ours: it is the one \cite{TY} applies to the Robertson graph, and is
attributed there to earlier work of Beveridge et al.

\begin{lemma}\label{lem:cornerfree}
If $G$ has girth at least $5$ and $\delta(G)\ge2$, then $G$ has no corner.
\end{lemma}

\begin{proof}
Suppose $N(x)\subseteq N[u]$ with $u\ne x$. Since $\deg x\ge2$ we may pick distinct
$a,b\in N(x)$. If $u\notin\{a,b\}$ then $a,b\in N(u)$ and $x\,a\,u\,b\,x$ is a $4$-cycle. If
$u=a$ then $b\in N(x)\cap N[a]$ and $b\ne a$, so $ab\in E(G)$ and $x\,a\,b\,x$ is a triangle.
Either way the girth is at most $4$.
\end{proof}

\begin{lemma}\label{lem:retract}
If $H$ is a retract of $G$ with $|V(H)|=|V(G)|-1$, then the vertex of
$V(G)\setminus V(H)$ is a corner of $G$. Consequently a corner-free graph has no retract of
order $|V(G)|-1$, and, since retractions compose, none is reached by a chain of retractions
either.
\end{lemma}

\begin{proof}
Let $V(G)\setminus V(H)=\{v\}$ and let $f$ witness the retraction. For $y\in N(v)$ we have
$y\in V(H)$, so $f(y)=y$, and $vy\in E(G)$ forces $f(v)y\in E(H)$ or $f(v)=y$; either way
$y\in N[f(v)]$. Hence $N(v)\subseteq N[f(v)]$ with $f(v)\ne v$, so $v$ is a corner. For the
last statement, a chain of retractions from order $|V(G)|$ to order $|V(G)|-1$ contains a step
that removes exactly one vertex, and a composite of retractions is a retraction.
\end{proof}

\begin{theorem}\label{thm:main}
$\cn(W_1)=\cn(W_2)=4$, and neither $W_1$ nor $W_2$ has any retract of order $19$. Since
$M_4=19$ and $|V(W_i)|=20=M_4+1$, the case $k=4$ of the question quoted in
Section~\ref{sec:question} is false.
\end{theorem}

\begin{proof}
Both graphs have girth $5$ and are $4$-regular, so $\cn(W_i)\ge\delta(W_i)=4$ by
Theorem~\ref{thm:af}. For the upper bound apply Lemma~\ref{lem:forest} with
\[
  S=\{0,14,16\}\ \text{in } W_1,
  \qquad
  S=\{0,2,5\}\ \text{in } W_2 .
\]
In $W_1$, $|N[S]|=14$ and $R=\{2,5,11,12,15,19\}$ induces exactly the five edges
$2\text{-}12$, $2\text{-}19$, $5\text{-}15$, $5\text{-}19$, $11\text{-}19$: a tree on six
vertices. In $W_2$, $|N[S]|=13$ and $R=\{1,3,4,7,10,12,19\}$ induces exactly the six edges
$1\text{-}7$, $1\text{-}12$, $3\text{-}10$, $4\text{-}10$, $4\text{-}12$, $7\text{-}19$,
which form the path $3\,10\,4\,12\,1\,7\,19$. Both remainders are forests, so
$\cn(W_i)\le4$ and therefore $\cn(W_i)=4$; that is, each $W_i$ is $4$-cop-win in the sense of
\cite{TY}. By Lemma~\ref{lem:cornerfree} neither graph has a corner, so by
Lemma~\ref{lem:retract} neither has a retract of order $19$ -- a fortiori none that is
$4$-cop-win. Since $M_4+1=20=|V(W_i)|$, the assertion of the question fails at $k=4$.
\end{proof}

\section{Scope}\label{sec:scope}

We do not classify the $4$-cop-win graphs on $20$ vertices: others may exist, of any girth,
and no minimality or uniqueness is claimed. $M_4=19$ is used, not verified; it is the corollary
of \cite{TY} that makes $20$ the relevant order. Only $k=4$ of the question is settled: $k=2$
is refuted in \cite{TY} itself, in the very paragraph quoted, by $C_5$ and $C_4$; $k=3$ is true
there by the classification immediately preceding that paragraph; and $k\ge5$ is untouched,
since $M_5$ is unknown and nothing here bounds it beyond the $M_5\ge20$ already implied by
$M_4=19$.

\section{Verification}\label{sec:verif}

The numerical claims above were re-derived by an accompanying program, \texttt{verify.py},
which uses the Python standard library only, in exact integer arithmetic and with no
randomness; a transcript of its run is in \texttt{verify.output.txt}, and it reports $51$
checks, all passing. From the graph6 strings and the adjacency tables of
Section~\ref{sec:graphs}, transcribed into it as literals, it decodes each string with its own
bit-unpacker, re-encodes it, and checks the decoded graph against the adjacency table printed
beside it, so the two encodings are cross-checked against each other; it recomputes the order,
the size, the degrees, connectivity, the triangle count, $\max_{a\ne b}|N(a)\cap N(b)|$ and the
girth; it finds $0$ corners over all $20\cdot19=380$ ordered pairs of each graph; it re-checks
the two forest certificates, and all $\binom{20}{3}=1140$ three-element sets; it computes
$|\mathrm{Aut}(W_1)|=96$ and $|\mathrm{Aut}(W_2)|=20$, the check on the identification with
Meringer's census entries; and it decides $\cn(W_1)=\cn(W_2)=4$ outright with a bitmask
least-fixed-point solver for the $k$-cop game, so that the value does not rest on
Theorem~\ref{thm:af} or Lemma~\ref{lem:forest}.

That solver is controlled on published cop numbers not our own: it returns $1$ on $P_5$, $2$ on
$C_4$, $2$ on $C_5$, $3$ on the Petersen graph and $4$ on the Robertson graph \cite{Robertson},
the last two read from the graph6 strings

\begin{verbatim}
Petersen  = IheA@GUAo
Robertson = R???C@_E?KI_I`bGCi?f?ocAoGAQO?
\end{verbatim}

\noindent
which the program also confirms to be cubic of order $10$ and girth $5$, and $4$-regular of
order $19$ and girth $5$, respectively. The program generates no graphs, so it does not
re-derive the census of \cite{Meringer}, and it does not verify $M_4=19$.

\begin{thebibliography}{9}

\bibitem{AF}
M.~Aigner and M.~Fromme,
\emph{A game of cops and robbers},
Discrete Appl. Math. \textbf{8} (1984), no.~1, 1--12.
\href{https://doi.org/10.1016/0166-218X(84)90073-8}{doi:10.1016/0166-218X(84)90073-8}.

\bibitem{Meringer}
M.~Meringer,
\emph{Fast generation of regular graphs and construction of cages},
J. Graph Theory \textbf{30} (1999), no.~2, 137--146.
The graphs $W_1$ and $W_2$ are Graph~2 and Graph~1 of the census file
\texttt{20\_4\_5.asc} distributed with that work.

\bibitem{Robertson}
N.~Robertson,
\emph{The smallest graph of girth $5$ and valency $4$},
Bull. Amer. Math. Soc. \textbf{70} (1964), 824--825.

\bibitem{TY}
J.~Turcotte and S.~Yvon,
\emph{$4$-cop-win graphs have at least $19$ vertices},
Discrete Appl. Math. \textbf{301} (2021), 74--98.
\href{https://doi.org/10.1016/j.dam.2021.05.012}{doi:10.1016/j.dam.2021.05.012};
\href{https://arxiv.org/abs/2006.02998v2}{arXiv:2006.02998v2}. The question refuted here is
quoted verbatim in Section~\ref{sec:question} from the e-print; it is an unnumbered question in
the discussion, distinct from the numbered Conjecture of that paper.

\end{thebibliography}

\end{document}
